\section{Conclusion}
\label{sec:conclusion}

We have presented the second generation of \Ziren{} as an end-to-end verifiable-computation system for MIPS32 execution. Its executor feeds opcode-specific arithmetic tables; its instruction frame and shared bus arguments connect those tables to program, state and memory; its Jagged--\WHIR{} layer reduces columns of different heights to one batched opening; and its recursion and verifier turn shard proofs into one deployable artifact. Trace area, interaction count, proof size and accelerator throughput therefore describe one implemented CPU--GPU proof system, not a synthesisable processor.

We have also separated the evidence for this design into explicit layers. Conformance suites and an executable \Lean{} model test the emulator, while mechanically extracted statements expose determinism obligations in the constraint system. The current artifact closes \ndetproved{} of \ndettotal{} such statements and leaves recursion outside the extraction. At the proof-system layer, the unique-decoding accounting applies to the checked-in schedule, while end-to-end knowledge extraction remains conditional on the per-node and extractor-composition premises of Theorem~\ref{thm:composition}. The cross-shard digest remains a separate structured-relation assumption.

The performance study has identified trace area, lookup interactions, layer order and shard size as distinct cost levers. The largest item it leaves on the table is the cross-shard digest: every word touched in a shard costs two rows of that unit, one quarter of the block's area, and either a digest of the per-shard delta or a larger shard would reduce it, the latter bounded by the memory the lookup circuit's first layer needs. Its measurements characterise the pre-correction implementation and one block-proving workload; they do not establish corrected-release performance or a same-hardware ranking against other systems. The next release should unify one runtime configuration, soundness report, key set, serialised proof and benchmark bundle after the Poseidon2 correction. Completing the composition proof, quantifying the digest assumption and closing the remaining trace-uniqueness premises are the corresponding proof priorities.
